\section{Termination for ordinary terminal fourfold pairs}\label{sec:counts}

We prove termination of ordinary flips for compact K\"ahler fourfold
pairs with effective real boundary and log discrepancy $a(E)>1$
for every exceptional place.  Two finite counts
replace projectivity of the ambient fourfold: analytic cycle classes
control contractions, and low-discrepancy places control flipped surfaces.
The bases of successive projective contractions may vary.
We first establish these counts, including their generalized-pair form
needed for extraction later in the proof.

\subsection{Loss of analytic cycle classes}

For a compact complex analytic space $V$ and an integer $k\geq 0$, put
\[
 \begin{split}
 \mathcal C_k(V)
 &:=\operatorname{span}_{\R}\{[S]\in H_{2k}(V,\R):
                  S\subset V\text{ irreducible compact analytic},\ \dim S=k\},\\
 c_k(V)&:=\dim_{\R}\mathcal C_k(V).
 \end{split}
\]
Fundamental classes here use the complex orientation on the regular loci.
Compact analytic spaces admit finite triangulations compatible with any
specified finite collection of closed analytic subsets, with simplicial
dimension at most twice the complex dimension.  Thus $c_k(V)$ is finite.
We justify this topological input after Lemma~\ref{lem:cycle-loss}.
For such a triangulation, Borel--Moore localization is the relative
simplicial-chain sequence for a closed analytic subset and its complement.

If $V$ is K\"ahler and $S$ is such a subspace, then
\begin{equation}\label{eq:positive-cycle}
 \langle[\omega]^k,[S]\rangle=\int_{S_{\mathrm{reg}}}\omega^k>0
\end{equation}
for a K\"ahler form $\omega$ on $V$.  Thus $[S]\neq 0$, including when
$V$ or $S$ is singular.

\begin{lemma}[Cycle loss]\label{lem:cycle-loss}
Let $X\dashrightarrow Y$ be a bimeromorphic transformation of compact
K\"ahler spaces, represented by a proper bimeromorphic diagram.
Suppose it identifies
\[
 U=X\setminus A\simeq Y\setminus B,
\]
where $A$ and $B$ are closed analytic subspaces and $\dim B<k$.
Then there is a surjection
\[
 \mathcal C_k(X)\longrightarrow\mathcal C_k(Y).
\]
If $A$ contains an irreducible compact analytic $k$-dimensional subspace,
then $c_k(X)>c_k(Y)$.

Consequently, in a sequence of bimeromorphic transformations of
$n$-dimensional compact K\"ahler spaces which extract no prime divisors,
only finitely many transformations contract a prime divisor.  A small
transformation preserves $c_{n-1}$.
\end{lemma}

\begin{proof}
The Borel--Moore localization sequence for $B\subset Y$ contains
\[
 H_{2k}(B,\R)\longrightarrow H_{2k}(Y,\R)
 \xrightarrow{\ j_Y^*\ }H_{2k}^{\mathrm{BM}}(U,\R)
 \longrightarrow H_{2k-1}(B,\R).
\]
The outside groups vanish: the real dimension of $B$ is at most $2k-2$.
Thus $j_Y^*$ is an isomorphism.  Compose restriction from $X$ with its
inverse to obtain a linear map
\[
 H_{2k}(X,\R)\xrightarrow{\ j_X^*\ }H_{2k}^{\mathrm{BM}}(U,\R)
 \xrightarrow{\ (j_Y^*)^{-1}\ }H_{2k}(Y,\R).
\]
For an irreducible $k$-cycle $S\subset X$ not contained in $A$, its
restriction is the fundamental class of $S\cap U$.  Closing its image in
$Y$ gives its strict transform: analyticity follows by taking the proper
image of the corresponding strict transform in the given diagram.
The displayed map therefore sends $[S]$ to this transformed cycle class.
It sends $[S]$ to zero when $S\subset A$.

Conversely, no $k$-dimensional subspace of $Y$ is contained in $B$.
Its strict transform in $X$ consequently maps to its fundamental class.
This proves the asserted surjection of cycle spans.  If $S\subset A$
has dimension $k$, its nonzero class, by \eqref{eq:positive-cycle}, is
in the kernel.  This proves the strict inequality.

For a transformation extracting no prime divisor, choose the common
isomorphic open so that the complement in the target has codimension
at least two.  Every contracted source prime lies in the source
complement.  Apply the result with $k=n-1$ at every step.  The
nonnegative integer $c_{n-1}$ cannot decrease infinitely often.  For a
small transformation both complements have codimension at least two,
so applying the result in both directions gives equality.
\end{proof}

This argument does not assert surjectivity on full homology, and does
not compare the Bott--Chern spaces of the two models.

\begin{proof}[Topological justification]
We give the abstract-space reduction for the finite compatible
triangulation used above.  Teissier's compatible Whitney stratification
\cite[Chapter~III, Propositions~1.5 and~2.2.2]{Teissier1982} has finitely
many strata on a compact space.  Choose finitely many analytic charts
$z_j:U_j\hookrightarrow\C^{N_j}$ and nonnegative smooth cutoffs
$\rho_j$ compactly supported in these charts whose positivity sets
cover the space.  Extend $\rho_j$ and $\rho_j z_j$ by zero.  The map
\[
 x\longmapsto\bigl(\rho_j(x),\rho_j(x)z_j(x)\bigr)_j
\]
is a topological embedding into a Euclidean space.  Near a point where
$\rho_j>0$, its components extend smoothly to the ambient chart and the
$j$-th block has the smooth left inverse $(t,w)\mapsto w/t$.
It is therefore locally an ambient smooth embedding, which preserves
the Whitney conditions.  Mather's control data
\cite[Section~7, Proposition~7.1, and Section~8]{Mather1970} and
Goresky's triangulation theorem \cite[Section~5]{Goresky1978} now give
a compatible triangulation, smooth on each stratum.  Compactness makes
it finite, and smoothness on strata gives the dimension bound.
Borel--Moore localization is then the relative simplicial-chain
sequence for a closed analytic subset and its complement.
\end{proof}

\subsection{Low places on compact log smooth data}

We use generalized discrepancies in the sense of
Section~\ref{sec:analytic}.  A place is called exceptional over $X$ when
its center on $X$ has codimension at least two.
The generalized klt statement below supplies the finite extraction set in
Section~\ref{sec:scaling}.  Its coordinate estimate and finite centre-blowup
argument also enter the terminal surface analysis in this section.

\begin{lemma}[Low places and their projective realization]\label{lem:low-places}
Let $(X,B+\bM)$ be a generalized pair on a compact normal complex
space, with effective real boundary of finite support and nef
b-$(1,1)$ data descending to a model projective over $X$.  Fix a projective log resolution $p:W\to X$ carrying those data.
\begin{enumerate}[label=\textup{(\roman*)}]
\item If the pair is generalized klt, there is $\epsilon>0$ such that
every prime divisor over $X$ has log discrepancy at least $\epsilon$,
and only finitely many exceptional places have
\[
 a(E;X,B+\bM)<1+\epsilon.
\]
\item If the pair is generalized terminal and
$b=\max(\{0\}\cup\{\operatorname{coeff}_D B\})$, there are only
finitely many exceptional places with
\[
 a(E;X,B+\bM)<2-b.
\]
\end{enumerate}
In each assertion the finite set of exceptional places can be represented
simultaneously by prime divisors on a smooth model projective over $X$.
In particular, any prescribed finite set of exceptional places with
$a(E;X,B+\bM)\leq 1$ for a generalized klt pair has such a realization.
This conclusion does not use an extraction or termination theorem.
\end{lemma}

\begin{proof}
Let $L\in\BC(X)$ be the generalized log class and let
$M_W=[\bM_W]$ be the nef trace class on $W$.  Write the crepant formula as
\begin{equation}\label{eq:carrier-boundary-count}
 [K_W+\Delta]+M_W=p^*L.
\end{equation}
This is the crepant boundary identity defining generalized
discrepancies; it is an equality of Bott--Chern classes when the nef
data are transcendental.  The boundary $\Delta$ is an actual real
divisor with simple normal crossing support.  Since the nef data
descend to $W$, every further generalized discrepancy is the ordinary
discrepancy for $(W,\Delta)$.  The coefficients of $\Delta$ need not be
nonnegative.  We may enlarge its listed simple normal crossing divisor
by components of coefficient zero.
Here a closed stratum means an irreducible component of an intersection
of listed components; it is smooth, and is the closure of its open
stratum.

We give the local estimate and the realization argument explicitly;
they are the log smooth calculations underlying
\cite[Proposition~2.36]{KM1998} and
\cite[Proposition~2.8]{CT2023}.
Write $\Delta=\sum_j d_jD_j$ and put $w_j=1-d_j$.
Suppose first that every $w_j$ is positive, and set
\[
 \delta=\min\bigl(\{1\}\cup\{w_j\}_j\bigr)>0.
\]
Let $E$ be exceptional over a smooth model with this log smooth
boundary.  At a general point of its center $C$, let $c=\codim C$,
and choose local parameters $x_1,\ldots,x_c$ for $C$ so that the
boundary components containing $C$ are $x_1=0,\ldots,x_r=0$.
For $v=\operatorname{ord}_E$, the Jacobian calculation gives
\begin{equation}\label{eq:coordinate-discrepancy-count}
 a(E;W,\Delta)\geq
 \sum_{j=1}^{r}(1-d_j)v(x_j)+\sum_{j=r+1}^{c}v(x_j).
\end{equation}
Indeed, on a smooth model containing $E$, with local parameter $t$
for $E$, write $x_j=t^{v(x_j)}u_j$ at its general point.  In a top
exterior differential at most one factor can use the term involving
$dt$ that lowers the order by one.  Thus the relative Jacobian has
order at least $\sum_{j=1}^c v(x_j)-1$.  Adding one and subtracting
the boundary orders proves \eqref{eq:coordinate-discrepancy-count}.
The remaining coordinates along $C$ are holomorphic and contribute
no negative orders.

In particular every discrepancy is at least $\delta$: for a divisor
already on the smooth model it is its listed weight, or one if it is
not listed; for an exceptional divisor use
\eqref{eq:coordinate-discrepancy-count}.  If $C$ is not a stratum of
the listed simple normal crossing divisor, then $c>r$.  When $r>0$,
the right hand side is at least $\delta+1$; when $r=0$, it is at
least $c\geq 2\geq\delta+1$.  Hence
\begin{equation}\label{eq:nonstratum-cutoff}
 C\text{ not a stratum}\quad\Longrightarrow\quad
 a(E;W,\Delta)\geq 1+\delta.
\end{equation}

Now fix a place $E$ with discrepancy strictly below $1+\delta$.
If it is not already a divisor on $W$, its center is a stratum by
\eqref{eq:nonstratum-cutoff}.  Blow up that closed stratum.
The center is smooth, the blowup is projective, and the new listed
boundary is again simple normal crossing.  If the center is the
intersection of the components indexed by $J$, then the new exceptional
component has weight
\begin{equation}\label{eq:stratum-weight-sum}
 w_{\mathrm{new}}=\sum_{j\in J}w_j.
\end{equation}
Thus all weights remain at least $\delta$, so the same argument
applies on the next model.  It excludes a non-stratum center at
every stage until $E$ becomes a divisor.

This process is finite.  While $E$ is exceptional, a blowup of its
smooth center of codimension $c\geq2$ changes its ordinary
empty-boundary discrepancy by
\[
 a(E;W_{\mathrm{new}},0)
 =a(E;W_{\mathrm{old}},0)-(c-1)v(\mathcal I_C).
\]
The subtracted quantity is a positive integer, whereas an ordinary
log discrepancy over a smooth space is a positive integer.
Consequently only finitely many such blowups are possible.
Thus $E$ is obtained by successive stratum blowups, which is the
meaning of a \emph{toroidal place} here.  This also realizes it on a projective model without first assuming
that the original model on which $E$ was given was projective over $X$.

For completeness, there are only finitely many such toroidal places.
A toroidal prime over a fixed irreducible stratum is determined by a
primitive nonzero vector $(m_j)_{j\in J}$ of nonnegative integers
in the normal boundary directions.  Its discrepancy is
\begin{equation}\label{eq:monomial-discrepancy-count}
 \sum_{j\in J}m_jw_j.
\end{equation}
These statements follow also directly from the successive stratum
blowups: a blowup inserts the sum of the relevant coordinate rays,
and the discrepancy transforms by \eqref{eq:stratum-weight-sum}.
The vector determines the monomial order and is unchanged when a
boundary equation is multiplied by a unit.  There are finitely many
irreducible strata on the compact model.  Since each $w_j\geq\delta$,
an upper bound $T$ on \eqref{eq:monomial-discrepancy-count} bounds
every $m_j$ by $T/\delta$.  This leaves finitely many vectors and
hence finitely many places.  Add the finitely many prime divisors
already on $W$ which are exceptional over $X$.

For a generalized klt pair, all weights in
\eqref{eq:carrier-boundary-count} are positive.  The preceding
argument with $\epsilon=\delta$ proves \textup{(i)} and its projective
realization assertion.  In particular $a\leq1$ lies strictly below
$1+\delta$.

No rationality of the weights was used: their fixed positive minimum
bounds the nonnegative integer vectors.  Thus the conclusion applies
to real boundaries and real nef data.

For a generalized terminal pair, the coefficient of every
$p$-exceptional component in $\Delta$ is negative, since its
log discrepancy is greater than one.  The other coefficients are
the original boundary coefficients and are at most $b<1$.
All the weights on $W$ are therefore at least $1-b$.
Repeat the argument with $\delta=1-b$ and $T=2-b$ to prove
\textup{(ii)}.  Subsequent stratum blowups preserve this lower bound
by \eqref{eq:stratum-weight-sum}.

Finally, to represent the entire finite set at once, take the
dominating component of the fiber product of the finitely many
projective realizations over $W$, normalize, and resolve projectively.
The strict transforms of the designated prime divisors remain
divisors on this common model.  Its composite map to $X$ is projective.
\end{proof}

\begin{remark}\label{rem:strict-low-cutoff}
The strict inequalities in Lemma~\ref{lem:low-places} are necessary.
For example, in $\mathbf P^n$, $n\geq3$, let $H$ be a hyperplane
with coefficient $b\in[0,1)$.  Blowing up any codimension-two linear
subspace contained in $H$ gives log discrepancy $2-b$ for
$(\mathbf P^n,bH)$.  There are infinitely many such places, although
this pair is terminal.  With empty boundary their discrepancy is two.
Below we count places strictly below a threshold and use a step at
which a place reaches that threshold to obtain a strict decrease.
\end{remark}


\subsection{Surface centres of low-discrepancy places}

The terminal termination argument must allow many places of discrepancy
below two above a moving boundary surface.  After recalling smoothness in
codimension two, Lemma~\ref{lem:terminal-low-centres} identifies the infinite
families that will be harmless.  Excluding only
the first blowup above each surface would be incorrect; the correction
in \cite[Proposition~3.1 and Lemma~3.2]{Fujino2005Addendum} is essential.

\begin{lemma}\label{lem:terminal-codim-two}
Let $(T,\Phi)$ be an ordinary klt pair on a normal complex fourfold,
with effective real boundary, $K_T$ a $\Q$-line bundle, and
$a(E;T,\Phi)>1$ for every exceptional prime divisor $E$ over $T$.
Then $T$ is smooth in codimension two.
\end{lemma}

\begin{proof}
Dropping the effective boundary shows that $T$ has terminal
singularities.  Suppose its singular locus has a codimension-two
component.  Near a general smooth point of this component choose a
holomorphic projection to $\C^2$ submersive along it.  Take a general
fibre transverse both to the regular locus and, on a fixed log
resolution, to the exceptional divisor and its strata.  These
transversality conditions hold after shrinking and choosing a general
value.  The resulting surface is a complete intersection in the
Cohen--Macaulay space $T$, since klt singularities are rational, and is
regular outside isolated points.  It is therefore normal.

Adjunction of a local Cartier pluricanonical power, first on the regular
locus and then by reflexivity, restricts the relative canonical equality
to the surface resolution.  All its exceptional discrepancy coefficients
are positive.  A normal surface germ with this property is smooth.
Indeed, negative definiteness of the exceptional intersection matrix
shows that a positive exceptional discrepancy divisor has negative
intersection with some exceptional curve.  That curve has negative
canonical degree; adjunction makes it a smooth rational $(-1)$-curve.
Contract it smoothly and repeat.  The remaining discrepancy coefficients
stay positive by pushforward.  Eventually no exceptional curve remains,
and the proper bimeromorphic map to the normal surface is an isomorphism.

On the other hand, a slice of a singular fourfold point by two equations
has embedding dimension at least the ambient embedding dimension minus
two, and hence greater than two.  The chosen surface is singular there,
a contradiction.
\end{proof}

\begin{lemma}\label{lem:terminal-low-centres}
Let $(T,\Phi)$ be an ordinary klt pair on a compact normal complex
fourfold, where $\Phi\geq0$ is real and
$a(E;T,\Phi)>1$ for every exceptional prime divisor $E$ over $T$.
Apart from finitely many exceptional places, every exceptional place with
$a(E;T,\Phi)<2$ has centre a surface $S\subset T$ with the following
properties: $T$ and the reduced support of $\Phi$ are generically smooth
along $S$, exactly one positive boundary component contains $S$, and,
if its coefficient is $b$, then
\[
 a(E;T,\Phi)\geq 2-b.
\]
All places above such a surface are allowed in this conclusion, not just
the exceptional divisor of its first blowup.
\end{lemma}

\begin{proof}
Choose a log resolution on which the positive strict boundary components
are smooth and mutually disjoint.  Starting with simple normal crossing
support, blow up intersections of the positive components to separate
them.  Every exceptional crepant coefficient is negative by terminality.
Thus the only positive components on this resolution are those disjoint
strict transforms, with their original coefficients.

Consider a place not already on the resolution.  At its general centre,
apply \eqref{eq:coordinate-discrepancy-count}, dropping negative boundary
coefficients.  A centre of codimension at least three has log discrepancy
at least $3-b>2$, if it lies in a positive component of coefficient $b$,
and at least three otherwise.  A codimension-two centre outside the
positive support has log discrepancy at least two.  Thus a place with
$a<2$ has codimension-two centre in one positive component, and its log
discrepancy is at least $2-b$.

If this centre lies in the exceptional locus, it is an irreducible
component of the intersection of that positive component and an
exceptional divisor.  There are finitely many such centres on the compact
resolution.  We check that each supports only finitely many places with
$a<2$.  It suffices to retain only the smooth component of coefficient
$b$, since deleting the negative coefficients lowers discrepancies.
Make this reduction once on the fixed resolution; on subsequent blowups
retain all exceptional coefficients in the crepant boundary.
Blow up the fixed centre at its general point.  Unless the given place
has appeared, the same estimate forces its next centre to be the
intersection of the positive strict transform and the new exceptional
divisor.  Indeed, the strict transform of the positive divisor maps isomorphically
to that divisor near the generic point of the original centre, because
the centre is Cartier within it. A codimension-two centre over that
generic point must therefore be this unique intersection. A proper
subset has codimension at least three and discrepancy at least $3-b$;
a centre outside the positive divisor has discrepancy at least two.
Thus no other branch is possible. The process reaches the given place
after finitely many blowups, by the decreasing positive integer
empty-boundary discrepancy used in Lemma~\ref{lem:low-places}.
For this reduced pair, the $k$th exceptional divisor on the branch has
crepant coefficient $-k(1-b)$ and log discrepancy
\[
 1+k(1-b),\qquad k=1,2,\ldots.
\]
Indeed, the first coefficient is $b-1$, and blowing up the intersection
of the newest exceptional divisor with the coefficient-$b$ strict
transform changes $-k(1-b)$ to $b-k(1-b)-1=-(k+1)(1-b)$.
Thus only
$k<1/(1-b)$ can occur before the cutoff two is reached.  This proves
finiteness over every fixed exceptional centre.  These generic blowups
are realized globally by blowing up the closures of the centres.

Add the finitely many exceptional primes on the chosen resolution.
All other centres under consideration lie outside its exceptional locus
at their general points, where the resolution is an isomorphism.  Their
images are precisely the surfaces described in the statement.
\end{proof}

We now have the two counts needed for termination.  A flipped surface
will give a place whose discrepancy strictly increases into a fixed
finite set.  Once such surfaces disappear, a compact surface-class rank
will decrease at every remaining step.

\subsection{Real terminal pairs}

\begin{proposition}\label{prop:terminal-real-four}
Let $(T_0,\Phi_0)$ be an ordinary klt pair on a normal globally
$\Q$-factorial compact K\"ahler fourfold, with $K_{T_0}$ a
$\Q$-line bundle and effective real boundary.  Suppose
$a(E;T_0,\Phi_0)>1$ for every exceptional prime divisor $E$ over $T_0$.
Every sequence of ordinary $(K_{T_i}+\Phi_i)$-flips starting at this
pair terminates, provided all models are normal globally $\Q$-factorial
compact K\"ahler fourfolds, each $K_{T_i}$ is a $\Q$-line bundle, and
each flip is given by projective small contractions
\[
 T_i\xrightarrow{\ f_i\ }Z_i
 \xleftarrow{\ f_i^+\ }T_{i+1}
\]
with $-(K_{T_i}+\Phi_i)$ relatively ample for $f_i$ and
$K_{T_{i+1}}+\Phi_{i+1}$ relatively ample for $f_i^+$.
Here $\Phi_{i+1}$ is the strict transform of $\Phi_i$.
No pseudo-effectivity or bigness assumption is required.
\end{proposition}

\begin{proof}
By Lemma~\ref{lem:negativity}, discrepancies do not decrease, and they
increase strictly for a place centred in either exceptional locus.
Since the maps are small, the exceptional places are the same on all
models.  Thus every pair remains terminal in the stated sense, and its
ambient fourfold is smooth in codimension two by
Lemma~\ref{lem:terminal-codim-two}.

Write the distinct positive coefficients of the boundary in decreasing
order.  We successively remove, by discarding finitely many initial
steps, every flipping and flipped surface contained in a component of
each coefficient.  After the positive coefficients we treat the value
zero, at which stage every flipped surface is considered.

\smallskip\noindent
\emph{Eliminating flipped surfaces.}
Fix the current coefficient $b$, and suppose that no flipping or flipped
surface is contained in a component of coefficient greater than $b$.
For $b=0$ assume that all positive coefficients have already been treated.
If $S$ is a flipped surface contained in a coefficient-$b$ component,
or any flipped surface when $b=0$, blow it up at its general point.
The exceptional prime $E$ has new log discrepancy
\begin{equation}\label{eq:surface-witness-values}
 a(E;T_{i+1},\Phi_{i+1})
   =2-\sum_{l=1}^{s}b_l\operatorname{mult}_{S}\Phi_{l,i+1},
 \qquad 1<a(E;T_{i+1},\Phi_{i+1})\leq2-b,
\end{equation}
where $\Phi_{l,i+1}$ are the positive prime boundary components and
$b_l$ their fixed coefficients.  The blowup place is defined globally
by blowing up the closed surface and resolving; its generic point
suffices for this computation.

The values in \eqref{eq:surface-witness-values} belong to a finite set,
even for irrational coefficients.  Indeed, terminality gives
\[
 \sum_l b_l\operatorname{mult}_{S}\Phi_{l,i+1}<1,
 \qquad
 0\leq\operatorname{mult}_{S}\Phi_{l,i+1}<1/b_l.
\]
Each multiplicity is an integer, and there are finitely many positive
components.  We use no discreteness assertion for all discrepancies.

Strictness in Lemma~\ref{lem:negativity} gives
$a(E;T_i,\Phi_i)<2-b$.  Its discrepancy on the first model of the
current tail is therefore also strictly below $2-b$.  By
Lemma~\ref{lem:terminal-low-centres}, outside a finite set such a place
would initially have a surface centre in one boundary component of
coefficient $c$, with discrepancy at least $2-c$.  Necessarily $c>b$.
That centre persists as a surface in this component at every later
step: at each step the map is an isomorphism at its general point,
since otherwise the centre itself would be a flipping surface in a
higher-coefficient component.  It cannot become a flipped surface
either, by the same higher-coefficient exclusion.  Thus it cannot be
the place just constructed.  All these witness places belong to a
fixed finite set.

Each member of that finite set is used only finitely often.  Between
uses its discrepancy does not decrease, and at each use it strictly
increases to one of the finitely many values
\eqref{eq:surface-witness-values}.  Consequently only finitely many
steps have a flipped surface of the current kind.  Pass beyond them.

\smallskip\noindent
\emph{Eliminating flipping surfaces inside the boundary.}
Suppose $b>0$.  Fix a coefficient-$b$ component and normalize its
transforms on the two sides of a flip.  Denote the normalizations by
$D_i^\nu,D_{i+1}^\nu$, and normalize their common image in $Z_i$ to
obtain $B_i$.  The induced maps to $B_i$ are proper and bimeromorphic.
They are isomorphisms away from the inverse image of the common
exceptional image in $Z_i$, whose dimension is at most one by
Lemma~\ref{lem:negativity}.  Finite normalization preserves this bound.
On $D_{i+1}^\nu$ its inverse image also has dimension at most one,
because no flipped surface is contained in the boundary component.

Every compact surface in $B_i$ has a strict transform on $D_i^\nu$,
so pushforward gives a surjection
\[
 \mathcal C_2(D_i^\nu)\longrightarrow\mathcal C_2(B_i).
\]
On the other side, Borel--Moore localization along the exceptional
subset of dimension at most one makes restriction in degree four
injective.  Compatibility with proper pushforward and the common
isomorphic open therefore gives an injection
\[
 \mathcal C_2(D_{i+1}^\nu)\lhook\joinrel\longrightarrow
 \mathcal C_2(B_i).
\]
These are assertions about surface spans; no surjectivity on the full
fourth homology of the negative side is needed. Only compact analytic
topology is used on the normalized components and their common image;
the positivity needed next comes from the ambient Kähler fourfold.

If a flipping surface lies in this boundary component, a surface above
it on $D_i^\nu$ maps to a set of dimension at most one in $B_i$.
Its class is nonzero: its image in $T_i$ is a surface, and the square
of a K\"ahler class on $T_i$ has strictly positive integral over it.
One may compute this integral on a resolution of the surface.
The finite normalization has positive generic degree on this surface,
so the same nonvanishing holds upstairs.  Its class is a nonzero
kernel vector in the displayed surjection.  Hence
\[
 c_2(D_i^\nu)\geq c_2(D_{i+1}^\nu),
\]
with strict inequality whenever the component contains a flipping
surface.  Summing over the finitely many coefficient-$b$ components
shows that only finitely many such steps occur.  This completes the
induction at $b$.

\smallskip\noindent
\emph{The remaining flips.}
After the positive coefficients have been treated, perform the
flipped-surface argument with $b=0$.  No flipped surface remains.
For every remaining flip, Lemma~\ref{lem:negativity} gives
\[
 \dim\Exc(f_i)+\dim\Exc(f_i^+)\geq3.
\]
Smallness bounds both dimensions by two, so $\Exc(f_i)$ contains a
surface and $\dim\Exc(f_i^+)\leq1$.  Apply
Lemma~\ref{lem:cycle-loss} with $k=2$ to the common isomorphic open.
It gives $c_2(T_i)>c_2(T_{i+1})$ at every remaining step, impossible
for an infinite sequence of nonnegative integers.
\end{proof}

The algebraic discrepancy and cycle-count strategy originates in
\cite[Theorem~5-1-15 and Lemmas~5-1-16--5-1-17]{KMM1987}
and \cite{Fujino2004,Fujino2005Addendum}.
The proof above supplies the real-boundary argument and uses compact
analytic cycle classes with K\"ahler positivity.  It applies those
geometric tools one step at a time, without a common projective base.

\begin{corollary}\label{prop:terminal-four}
Let $X_0$ be a normal globally strongly $\Q$-factorial compact
K\"ahler fourfold with terminal singularities.
Every sequence starting at $X_0$ of ordinary empty-boundary negative
divisorial contractions and flips, with the projectivity and compact
K\"ahler conclusions of Proposition~\ref{prop:ordinary-step}, terminates.
No pseudo-effectivity hypothesis is required for this assertion about
birational sequences.
\end{corollary}

\begin{proof}
For flips, terminality persists by Lemma~\ref{lem:negativity}.
For a divisorial contraction, \eqref{eq:ordinary-divisorial-difference}
gives $K_{X_i}=f^*K_{X_{i+1}}+eE$ with $e>0$.
The lost prime has target log discrepancy $1+e>1$; every other
exceptional place has target discrepancy at least its source
discrepancy, which was greater than one.  Thus terminality persists
throughout the sequence.
Lemma~\ref{lem:cycle-loss} with $k=3$ bounds the number of divisorial
steps.  Any infinite sequence would therefore have a tail consisting
only of flips, contrary to Proposition~\ref{prop:terminal-real-four}
with zero boundary.
\end{proof}
